\documentclass{modern-report}

% =============================================================================
% CONFIGURATION DU RAPPORT - À PERSONNALISER
% =============================================================================

% Informations du rapport
\reporttitle{[TITRE DE VOTRE RAPPORT]}
\reportsubtitle{[Sous-titre descriptif]}
\reportauthor{[Votre Nom]}
\reportdate{\today}  % Ou une date spécifique : {15 janvier 2025}
\reportversion{Version 1.0}
\reportclient{[Nom du Client]}
\reportorganization{[Votre Organisation]}

\begin{document}

% =============================================================================
% PAGE DE GARDE MODERNE
% =============================================================================
\makemoderncover

% =============================================================================
% TABLE DES MATIÈRES
% =============================================================================
\tableofcontents
\clearpage

% =============================================================================
% RÉSUMÉ EXÉCUTIF
% =============================================================================
\chapter*{Résumé Exécutif}
\addcontentsline{toc}{chapter}{Résumé Exécutif}

\begin{modernbox}{Synthèse}
% Remplacez ce texte par votre résumé exécutif
Voici un résumé concis des points clés de votre rapport. Cette section doit présenter les éléments essentiels de manière claire et synthétique.

\modernseparator

\textbf{Objectifs principaux :}
\begin{itemize}
    \item Objectif 1
    \item Objectif 2 
    \item Objectif 3
\end{itemize}
\end{modernbox}

% =============================================================================
% CHAPITRES PRINCIPAUX
% =============================================================================

\chapter{Introduction}

\section{Contexte et Enjeux}

% Décrivez ici le contexte de votre projet ou étude
Présentez ici le contexte général de votre rapport, les enjeux identifiés et la problématique à résoudre.

\subsection{Problématique}

% Définissez la problématique centrale
Énoncez clairement la problématique que votre rapport traite.

\begin{infobox}{Information Importante}
Utilisez cette boîte pour mettre en évidence des informations importantes ou des points clés à retenir.
\end{infobox}

\section{Objectifs}

% Listez vos objectifs
Les objectifs de ce rapport sont les suivants :

\begin{enumerate}
    \item Premier objectif spécifique
    \item Deuxième objectif mesurable
    \item Troisième objectif atteignable
\end{enumerate}

\modernseparator

\chapter{Analyse et Méthodologie}

\section{Approche Méthodologique}

% Décrivez votre méthodologie
Présentez ici la méthodologie utilisée pour mener votre analyse ou réaliser votre projet.

\subsection{Phase 1 : Analyse Préliminaire}

% Décrivez la première phase
Détaillez les étapes de votre phase d'analyse préliminaire.

\subsection{Phase 2 : Mise en Œuvre}

% Décrivez la phase de mise en œuvre
Expliquez comment vous avez procédé pour la mise en œuvre de votre solution.

\begin{warningbox}{Point d'Attention}
Utilisez cette boîte pour signaler des points d'attention, des risques identifiés ou des précautions à prendre.
\end{warningbox}

\section{Outils et Technologies}

% Si applicable, présentez vos outils
\begin{table}[h]
\centering
\caption{Outils Utilisés}
\begin{tabular}{@{}ll@{}}
\toprule
\moderntablehead{Catégorie & Outil/Technologie} \\
\midrule
Analyse & [Nom de l'outil] \\
Développement & [Framework utilisé] \\
Base de données & [SGBD] \\
Interface & [Framework UI] \\
\bottomrule
\end{tabular}
\end{table}

\chapter{Résultats et Réalisations}

\section{Principales Réalisations}

% Présentez vos résultats
Détaillez ici vos principales réalisations et les résultats obtenus.

\subsection{Réalisation 1}

% Première réalisation
Description de votre première réalisation majeure.

\subsection{Réalisation 2}

% Deuxième réalisation
Description de votre deuxième réalisation importante.

\modernseparator

\section{Indicateurs de Performance}

% Si applicable, présentez vos KPI
\begin{modernbox}{Métriques Clés}
\begin{itemize}
    \item Indicateur 1 : [Valeur et amélioration]
    \item Indicateur 2 : [Résultat obtenu]
    \item Indicateur 3 : [Performance mesurée]
\end{itemize}
\end{modernbox}

\chapter{Recommandations et Perspectives}

\section{Recommandations}

% Vos recommandations
Sur la base de cette analyse, nous formulons les recommandations suivantes :

\begin{enumerate}
    \item \textbf{Recommandation 1} : Description détaillée de la première recommandation
    \item \textbf{Recommandation 2} : Explication de la deuxième recommandation
    \item \textbf{Recommandation 3} : Présentation de la troisième recommandation
\end{enumerate}

\section{Plan d'Action}

% Votre plan d'action
\begin{infobox}{Prochaines Étapes}
\begin{itemize}
    \item \textbf{Court terme} : Actions à mener dans les 3 mois
    \item \textbf{Moyen terme} : Objectifs à 6-12 mois
    \item \textbf{Long terme} : Vision stratégique à 2-3 ans
\end{itemize}
\end{infobox}

% =============================================================================
% CONCLUSION
% =============================================================================
\chapter*{Conclusion}
\addcontentsline{toc}{chapter}{Conclusion}

% Votre conclusion
Rédigez ici votre conclusion en reprenant les éléments clés de votre rapport et en soulignant les apports de votre travail.

\begin{modernbox}{Synthèse Finale}
\begin{itemize}
    \item Résultat principal obtenu
    \item Impact attendu ou mesuré
    \item Valeur ajoutée apportée
\end{itemize}
\end{modernbox}

\modernseparator

\textit{[Ajoutez ici vos remerciements ou informations de contact si nécessaire]}

% =============================================================================
% ANNEXES (OPTIONNEL)
% =============================================================================
\appendix

\chapter{Spécifications Techniques}
% Annexes techniques si nécessaire

\section{Configuration Système}

Si votre rapport traite d'aspects techniques, vous pouvez détailler ici les spécifications système requises.

\chapter{Documentation Complémentaire}
% Autres annexes

\section{Références}

% Listez vos références bibliographiques
\begin{itemize}
    \item Référence 1
    \item Référence 2
    \item Référence 3
\end{itemize}

\section{Glossaire}

% Définitions des termes techniques
\begin{description}
    \item[Terme 1] : Définition du premier terme
    \item[Terme 2] : Explication du deuxième terme
    \item[Terme 3] : Description du troisième terme
\end{description}

% =============================================================================
% FIN DU DOCUMENT
% =============================================================================

\end{document}